% math_commands.tex
% Common mathematical notation and commands for ML/AI papers
% Include with: % math_commands.tex
% Common mathematical notation and commands for ML/AI papers
% Include with: % math_commands.tex
% Common mathematical notation and commands for ML/AI papers
% Include with: % math_commands.tex
% Common mathematical notation and commands for ML/AI papers
% Include with: \input{math_commands.tex}

% =============================================================================
% General Notation
% =============================================================================

% Vectors and matrices
\newcommand{\vect}[1]{\boldsymbol{\mathbf{#1}}}
\newcommand{\mat}[1]{\boldsymbol{\mathbf{#1}}}
\newcommand{\tensor}[1]{\mathcal{#1}}

% Transpose
\newcommand{\T}{^{\top}}
\newcommand{\Tr}{^{\intercal}}

% Inverse
\newcommand{\inv}{^{-1}}

% Hermitian conjugate
\newcommand{\herm}{^{\dagger}}

% Vector and matrix norms
\newcommand{\norm}[1]{\left\|#1\right\|}
\newcommand{\normp}[2]{\left\|#1\right\|_{#2}}
\newcommand{\abs}[1]{\left|#1\right|}

% Trace and determinant
\DeclareMathOperator{\tr}{tr}
\DeclareMathOperator{\Tr}{Tr}
\DeclareMathOperator{\diag}{diag}
\DeclareMathOperator{\Det}{det}

% Sets
\newcommand{\R}{\mathbb{R}}
\newcommand{\N}{\mathbb{N}}
\newcommand{\Z}{\mathbb{Z}}
\newcommand{\C}{\mathbb{C}}

% =============================================================================
% Probability and Statistics
% =============================================================================

% Probability operators
\DeclareMathOperator{\E}{\mathbb{E}}
\DeclareMathOperator{\Var}{Var}
\DeclareMathOperator{\Cov}{Cov}
\DeclareMathOperator{\KL}{KL}
\DeclareMathOperator{\Entropy}{H}
\DeclareMathOperator{\MI}{I}

% Probability commands
\newcommand{\prob}[1]{\mathbb{P}\left(#1\right)}
\newcommand{\expect}[1]{\E\left[#1\right]}
\newcommand{\expectsub}[2]{\E_{#1}\left[#2\right]}
\newcommand{\var}[1]{\Var\left[#1\right]}
\newcommand{\cov}[1]{\Cov\left[#1\right]}

% Distributions
\DeclareMathOperator{\Ndist}{\mathcal{N}}
\DeclareMathOperator{\Unif}{Uniform}
\DeclareMathOperator{\Bern}{Bernoulli}
\DeclareMathOperator{\Cat}{Categorical}

% Gaussian
\newcommand{\gauss}[2]{\Ndist\left(#1, #2\right)}

% =============================================================================
% Optimization
% =============================================================================

% Argmin/argmax
\DeclareMathOperator{\argmin}{arg\,min}
\DeclareMathOperator{\argmax}{arg\,max}
\DeclareMathOperator{\arginf}{arg\,inf}
\DeclareMathOperator{\argsup}{arg\,sup}

% Min/max
\DeclareMathOperator{\minimize}{minimize}
\DeclareMathOperator{\maximize}{maximize}

% Subgradient
\DeclareMathOperator{\sign}{sign}
\DeclareMathOperator{\prox}{prox}

% =============================================================================
% Linear Algebra
% =============================================================================

% Inner products
\newcommand{\inner}[2]{\left\langle #1, #2 \right\rangle}
\newcommand{\dotp}[2]{#1 \cdot #2}

% Outer product
\newcommand{\outerp}[2]{#1 \otimes #2}

% Matrix operations
\DeclareMathOperator{\rank}{rank}
\DeclareMathOperator{\nullity}{nullity}
\DeclareMathOperator{\range}{range}
\DeclareMathOperator{\nullspace}{null}
\DeclareMathOperator{\spanop}{span}

% Eigenvalues
\DeclareMathOperator{\eig}{eig}
\DeclareMathOperator{\eigmax}{\lambda_{max}}
\DeclareMathOperator{\eigmin}{\lambda_{min}}

% Singular values
\DeclareMathOperator{\svd}{svd}
\DeclareMathOperator{\sigmamax}{\sigma_{max}}
\DeclareMathOperator{\sigmamin}{\sigma_{min}}

% =============================================================================
% Machine Learning Notation
% =============================================================================

% Dataset
\newcommand{\dataset}{\mathcal{D}}
\newcommand{\train}{\mathcal{D}_{\text{train}}}
\newcommand{\test}{\mathcal{D}_{\text{test}}}
\newcommand{\val}{\mathcal{D}_{\text{val}}}

% Data samples
\newcommand{\x}{\vect{x}}
\newcommand{\y}{y}
\newcommand{\X}{\mat{X}}
\newcommand{\Y}{\vect{y}}

% Labels and predictions
\newcommand{\yhat}{\hat{y}}
\newcommand{\Yhat}{\hat{\vect{y}}}
\newcommand{\fhat}{\hat{f}}

% Parameters
\newcommand{\w}{\vect{w}}
\newcommand{\what}{\hat{\vect{w}}}
\newcommand{\wstar}{\vect{w}^{*}}
\newcommand{\theta}{\boldsymbol{\theta}}
\newcommand{\thetahat}{\hat{\boldsymbol{\theta}}}
\newcommand{\thetastar}{\boldsymbol{\theta}^{*}}

% Gradients
\newcommand{\grad}[1]{\nabla_{#1}}
\newcommand{\gradw}{\nabla_{\w}}
\newcommand{\gradtheta}{\nabla_{\theta}}

% Loss functions
\newcommand{\loss}[1]{\mathcal{L}\left(#1\right)}
\newcommand{\losshat}[1]{\hat{\mathcal{L}}\left(#1\right)}

% Regularization
\newcommand{\reg}[1]{\mathcal{R}\left(#1\right)}

% Objective
\newcommand{\obj}[1]{\mathcal{J}\left(#1\right)}

% Learning rate
\newcommand{\lr}{\eta}
\newcommand{\stepsize}{\alpha}

% Activation functions
\DeclareMathOperator{\relu}{ReLU}
\DeclareMathOperator{\sigmoid}{sigmoid}
\DeclareMathOperator{\softmax}{softmax}
\DeclareMathOperator{\tanhop}{tanh}

% Layer notation
\newcommand{\layer}[1]{^{[#1]}}

% =============================================================================
% Deep Learning
% =============================================================================

% Network
\newcommand{\net}{f_{\theta}}
\newcommand{\netn}{f_{\theta}^{(n)}}

% Attention
\DeclareMathOperator{\Attention}{Attention}
\DeclareMathOperator{\MultiHead}{MultiHead}
\DeclareMathOperator{\SelfAttention}{SelfAttention}

% Transformer
\DeclareMathOperator{\LayerNorm}{LayerNorm}
\DeclareMathOperator{\FFN}{FFN}

% Convolution
\DeclareMathOperator{\Conv}{Conv}
\DeclareMathOperator{\Pool}{Pool}

% =============================================================================
% Reinforcement Learning
% =============================================================================

% States and actions
\newcommand{\state}{s}
\newcommand{\action}{a}
\newcommand{\State}{\mathcal{S}}
\newcommand{\Action}{\mathcal{A}}

% Rewards
\newcommand{\reward}{r}
\newcommand{\Reward}{\mathcal{R}}

% Policy
\newcommand{\policy}{\pi}
\newcommand{\policytheta}{\pi_{\theta}}

% Value functions
\newcommand{\vval}{V}
\newcommand{\qval}{Q}
\newcommand{\aval}{A}

% Advantage
\newcommand{\advantage}{A}

% Return
\newcommand{\return}{G}

% Discount factor
\newcommand{\discount}{\gamma}

% =============================================================================
% Common Functions and Operators
% =============================================================================

% Big-O notation
\DeclareMathOperator{\bigO}{\mathcal{O}}
\DeclareMathOperator{\littleo}{o}
\DeclareMathOperator{\bigOmega}{\Omega}
\DeclareMathOperator{\bigTheta}{\Theta}

% Standard functions
\DeclareMathOperator{\softmaxop}{softmax}
\DeclareMathOperator{\logsumexp}{log-sum-exp}

% Indicator function
\newcommand{\indicator}[1]{\mathbb{I}\left[#1\right]}

% Cardinality
\newcommand{\card}[1]{\left|#1\right|}

% Floor and ceiling
\DeclareMathOperator{\floor}{floor}
\DeclareMathOperator{\ceil}{ceil}

% =============================================================================
% Brackets and Delimiters
% =============================================================================

% Parentheses
\newcommand{\paren}[1]{\left(#1\right)}
\newcommand{\bracket}[1]{\left[#1\right]}
\newcommand{\brace}[1]{\left\{#1\right\}}
\newcommand{\absval}[1]{\left|#1\right|}

% Set notation
\newcommand{\set}[1]{\left\{#1\right\}}

% Conditional
\newcommand{\given}{\,|\,}
\newcommand{\suchthat}{\,:\,}

% =============================================================================
% Common Abbreviations
% =============================================================================

% i.i.d.
\newcommand{\iid}{\textit{i.i.d.}}

% w.r.t., w.p., etc.
\newcommand{\wrt}{w.r.t.}
\newcommand{\wp}{w.p.}
\newcommand{\st}{s.t.}

% et al.
\newcommand{\etal}{\textit{et al.}}

% e.g., i.e., cf.
\newcommand{\eg}{e.g.,}
\newcommand{\ie}{i.e.,}
\newcommand{\cf}{cf.}

% =============================================================================
% Custom Operators (add your own below)
% =============================================================================

% Add your custom mathematical notation here
%
% \newcommand{\yourcommand}[n]{definition}
%
% Example:
% \newcommand{\model}{\mathcal{M}}
% \newcommand{\embed}[1]{\text{embed}(#1)}


% =============================================================================
% General Notation
% =============================================================================

% Vectors and matrices
\newcommand{\vect}[1]{\boldsymbol{\mathbf{#1}}}
\newcommand{\mat}[1]{\boldsymbol{\mathbf{#1}}}
\newcommand{\tensor}[1]{\mathcal{#1}}

% Transpose
\newcommand{\T}{^{\top}}
\newcommand{\Tr}{^{\intercal}}

% Inverse
\newcommand{\inv}{^{-1}}

% Hermitian conjugate
\newcommand{\herm}{^{\dagger}}

% Vector and matrix norms
\newcommand{\norm}[1]{\left\|#1\right\|}
\newcommand{\normp}[2]{\left\|#1\right\|_{#2}}
\newcommand{\abs}[1]{\left|#1\right|}

% Trace and determinant
\DeclareMathOperator{\tr}{tr}
\DeclareMathOperator{\Tr}{Tr}
\DeclareMathOperator{\diag}{diag}
\DeclareMathOperator{\Det}{det}

% Sets
\newcommand{\R}{\mathbb{R}}
\newcommand{\N}{\mathbb{N}}
\newcommand{\Z}{\mathbb{Z}}
\newcommand{\C}{\mathbb{C}}

% =============================================================================
% Probability and Statistics
% =============================================================================

% Probability operators
\DeclareMathOperator{\E}{\mathbb{E}}
\DeclareMathOperator{\Var}{Var}
\DeclareMathOperator{\Cov}{Cov}
\DeclareMathOperator{\KL}{KL}
\DeclareMathOperator{\Entropy}{H}
\DeclareMathOperator{\MI}{I}

% Probability commands
\newcommand{\prob}[1]{\mathbb{P}\left(#1\right)}
\newcommand{\expect}[1]{\E\left[#1\right]}
\newcommand{\expectsub}[2]{\E_{#1}\left[#2\right]}
\newcommand{\var}[1]{\Var\left[#1\right]}
\newcommand{\cov}[1]{\Cov\left[#1\right]}

% Distributions
\DeclareMathOperator{\Ndist}{\mathcal{N}}
\DeclareMathOperator{\Unif}{Uniform}
\DeclareMathOperator{\Bern}{Bernoulli}
\DeclareMathOperator{\Cat}{Categorical}

% Gaussian
\newcommand{\gauss}[2]{\Ndist\left(#1, #2\right)}

% =============================================================================
% Optimization
% =============================================================================

% Argmin/argmax
\DeclareMathOperator{\argmin}{arg\,min}
\DeclareMathOperator{\argmax}{arg\,max}
\DeclareMathOperator{\arginf}{arg\,inf}
\DeclareMathOperator{\argsup}{arg\,sup}

% Min/max
\DeclareMathOperator{\minimize}{minimize}
\DeclareMathOperator{\maximize}{maximize}

% Subgradient
\DeclareMathOperator{\sign}{sign}
\DeclareMathOperator{\prox}{prox}

% =============================================================================
% Linear Algebra
% =============================================================================

% Inner products
\newcommand{\inner}[2]{\left\langle #1, #2 \right\rangle}
\newcommand{\dotp}[2]{#1 \cdot #2}

% Outer product
\newcommand{\outerp}[2]{#1 \otimes #2}

% Matrix operations
\DeclareMathOperator{\rank}{rank}
\DeclareMathOperator{\nullity}{nullity}
\DeclareMathOperator{\range}{range}
\DeclareMathOperator{\nullspace}{null}
\DeclareMathOperator{\spanop}{span}

% Eigenvalues
\DeclareMathOperator{\eig}{eig}
\DeclareMathOperator{\eigmax}{\lambda_{max}}
\DeclareMathOperator{\eigmin}{\lambda_{min}}

% Singular values
\DeclareMathOperator{\svd}{svd}
\DeclareMathOperator{\sigmamax}{\sigma_{max}}
\DeclareMathOperator{\sigmamin}{\sigma_{min}}

% =============================================================================
% Machine Learning Notation
% =============================================================================

% Dataset
\newcommand{\dataset}{\mathcal{D}}
\newcommand{\train}{\mathcal{D}_{\text{train}}}
\newcommand{\test}{\mathcal{D}_{\text{test}}}
\newcommand{\val}{\mathcal{D}_{\text{val}}}

% Data samples
\newcommand{\x}{\vect{x}}
\newcommand{\y}{y}
\newcommand{\X}{\mat{X}}
\newcommand{\Y}{\vect{y}}

% Labels and predictions
\newcommand{\yhat}{\hat{y}}
\newcommand{\Yhat}{\hat{\vect{y}}}
\newcommand{\fhat}{\hat{f}}

% Parameters
\newcommand{\w}{\vect{w}}
\newcommand{\what}{\hat{\vect{w}}}
\newcommand{\wstar}{\vect{w}^{*}}
\newcommand{\theta}{\boldsymbol{\theta}}
\newcommand{\thetahat}{\hat{\boldsymbol{\theta}}}
\newcommand{\thetastar}{\boldsymbol{\theta}^{*}}

% Gradients
\newcommand{\grad}[1]{\nabla_{#1}}
\newcommand{\gradw}{\nabla_{\w}}
\newcommand{\gradtheta}{\nabla_{\theta}}

% Loss functions
\newcommand{\loss}[1]{\mathcal{L}\left(#1\right)}
\newcommand{\losshat}[1]{\hat{\mathcal{L}}\left(#1\right)}

% Regularization
\newcommand{\reg}[1]{\mathcal{R}\left(#1\right)}

% Objective
\newcommand{\obj}[1]{\mathcal{J}\left(#1\right)}

% Learning rate
\newcommand{\lr}{\eta}
\newcommand{\stepsize}{\alpha}

% Activation functions
\DeclareMathOperator{\relu}{ReLU}
\DeclareMathOperator{\sigmoid}{sigmoid}
\DeclareMathOperator{\softmax}{softmax}
\DeclareMathOperator{\tanhop}{tanh}

% Layer notation
\newcommand{\layer}[1]{^{[#1]}}

% =============================================================================
% Deep Learning
% =============================================================================

% Network
\newcommand{\net}{f_{\theta}}
\newcommand{\netn}{f_{\theta}^{(n)}}

% Attention
\DeclareMathOperator{\Attention}{Attention}
\DeclareMathOperator{\MultiHead}{MultiHead}
\DeclareMathOperator{\SelfAttention}{SelfAttention}

% Transformer
\DeclareMathOperator{\LayerNorm}{LayerNorm}
\DeclareMathOperator{\FFN}{FFN}

% Convolution
\DeclareMathOperator{\Conv}{Conv}
\DeclareMathOperator{\Pool}{Pool}

% =============================================================================
% Reinforcement Learning
% =============================================================================

% States and actions
\newcommand{\state}{s}
\newcommand{\action}{a}
\newcommand{\State}{\mathcal{S}}
\newcommand{\Action}{\mathcal{A}}

% Rewards
\newcommand{\reward}{r}
\newcommand{\Reward}{\mathcal{R}}

% Policy
\newcommand{\policy}{\pi}
\newcommand{\policytheta}{\pi_{\theta}}

% Value functions
\newcommand{\vval}{V}
\newcommand{\qval}{Q}
\newcommand{\aval}{A}

% Advantage
\newcommand{\advantage}{A}

% Return
\newcommand{\return}{G}

% Discount factor
\newcommand{\discount}{\gamma}

% =============================================================================
% Common Functions and Operators
% =============================================================================

% Big-O notation
\DeclareMathOperator{\bigO}{\mathcal{O}}
\DeclareMathOperator{\littleo}{o}
\DeclareMathOperator{\bigOmega}{\Omega}
\DeclareMathOperator{\bigTheta}{\Theta}

% Standard functions
\DeclareMathOperator{\softmaxop}{softmax}
\DeclareMathOperator{\logsumexp}{log-sum-exp}

% Indicator function
\newcommand{\indicator}[1]{\mathbb{I}\left[#1\right]}

% Cardinality
\newcommand{\card}[1]{\left|#1\right|}

% Floor and ceiling
\DeclareMathOperator{\floor}{floor}
\DeclareMathOperator{\ceil}{ceil}

% =============================================================================
% Brackets and Delimiters
% =============================================================================

% Parentheses
\newcommand{\paren}[1]{\left(#1\right)}
\newcommand{\bracket}[1]{\left[#1\right]}
\newcommand{\brace}[1]{\left\{#1\right\}}
\newcommand{\absval}[1]{\left|#1\right|}

% Set notation
\newcommand{\set}[1]{\left\{#1\right\}}

% Conditional
\newcommand{\given}{\,|\,}
\newcommand{\suchthat}{\,:\,}

% =============================================================================
% Common Abbreviations
% =============================================================================

% i.i.d.
\newcommand{\iid}{\textit{i.i.d.}}

% w.r.t., w.p., etc.
\newcommand{\wrt}{w.r.t.}
\newcommand{\wp}{w.p.}
\newcommand{\st}{s.t.}

% et al.
\newcommand{\etal}{\textit{et al.}}

% e.g., i.e., cf.
\newcommand{\eg}{e.g.,}
\newcommand{\ie}{i.e.,}
\newcommand{\cf}{cf.}

% =============================================================================
% Custom Operators (add your own below)
% =============================================================================

% Add your custom mathematical notation here
%
% \newcommand{\yourcommand}[n]{definition}
%
% Example:
% \newcommand{\model}{\mathcal{M}}
% \newcommand{\embed}[1]{\text{embed}(#1)}


% =============================================================================
% General Notation
% =============================================================================

% Vectors and matrices
\newcommand{\vect}[1]{\boldsymbol{\mathbf{#1}}}
\newcommand{\mat}[1]{\boldsymbol{\mathbf{#1}}}
\newcommand{\tensor}[1]{\mathcal{#1}}

% Transpose
\newcommand{\T}{^{\top}}
\newcommand{\Tr}{^{\intercal}}

% Inverse
\newcommand{\inv}{^{-1}}

% Hermitian conjugate
\newcommand{\herm}{^{\dagger}}

% Vector and matrix norms
\newcommand{\norm}[1]{\left\|#1\right\|}
\newcommand{\normp}[2]{\left\|#1\right\|_{#2}}
\newcommand{\abs}[1]{\left|#1\right|}

% Trace and determinant
\DeclareMathOperator{\tr}{tr}
\DeclareMathOperator{\Tr}{Tr}
\DeclareMathOperator{\diag}{diag}
\DeclareMathOperator{\Det}{det}

% Sets
\newcommand{\R}{\mathbb{R}}
\newcommand{\N}{\mathbb{N}}
\newcommand{\Z}{\mathbb{Z}}
\newcommand{\C}{\mathbb{C}}

% =============================================================================
% Probability and Statistics
% =============================================================================

% Probability operators
\DeclareMathOperator{\E}{\mathbb{E}}
\DeclareMathOperator{\Var}{Var}
\DeclareMathOperator{\Cov}{Cov}
\DeclareMathOperator{\KL}{KL}
\DeclareMathOperator{\Entropy}{H}
\DeclareMathOperator{\MI}{I}

% Probability commands
\newcommand{\prob}[1]{\mathbb{P}\left(#1\right)}
\newcommand{\expect}[1]{\E\left[#1\right]}
\newcommand{\expectsub}[2]{\E_{#1}\left[#2\right]}
\newcommand{\var}[1]{\Var\left[#1\right]}
\newcommand{\cov}[1]{\Cov\left[#1\right]}

% Distributions
\DeclareMathOperator{\Ndist}{\mathcal{N}}
\DeclareMathOperator{\Unif}{Uniform}
\DeclareMathOperator{\Bern}{Bernoulli}
\DeclareMathOperator{\Cat}{Categorical}

% Gaussian
\newcommand{\gauss}[2]{\Ndist\left(#1, #2\right)}

% =============================================================================
% Optimization
% =============================================================================

% Argmin/argmax
\DeclareMathOperator{\argmin}{arg\,min}
\DeclareMathOperator{\argmax}{arg\,max}
\DeclareMathOperator{\arginf}{arg\,inf}
\DeclareMathOperator{\argsup}{arg\,sup}

% Min/max
\DeclareMathOperator{\minimize}{minimize}
\DeclareMathOperator{\maximize}{maximize}

% Subgradient
\DeclareMathOperator{\sign}{sign}
\DeclareMathOperator{\prox}{prox}

% =============================================================================
% Linear Algebra
% =============================================================================

% Inner products
\newcommand{\inner}[2]{\left\langle #1, #2 \right\rangle}
\newcommand{\dotp}[2]{#1 \cdot #2}

% Outer product
\newcommand{\outerp}[2]{#1 \otimes #2}

% Matrix operations
\DeclareMathOperator{\rank}{rank}
\DeclareMathOperator{\nullity}{nullity}
\DeclareMathOperator{\range}{range}
\DeclareMathOperator{\nullspace}{null}
\DeclareMathOperator{\spanop}{span}

% Eigenvalues
\DeclareMathOperator{\eig}{eig}
\DeclareMathOperator{\eigmax}{\lambda_{max}}
\DeclareMathOperator{\eigmin}{\lambda_{min}}

% Singular values
\DeclareMathOperator{\svd}{svd}
\DeclareMathOperator{\sigmamax}{\sigma_{max}}
\DeclareMathOperator{\sigmamin}{\sigma_{min}}

% =============================================================================
% Machine Learning Notation
% =============================================================================

% Dataset
\newcommand{\dataset}{\mathcal{D}}
\newcommand{\train}{\mathcal{D}_{\text{train}}}
\newcommand{\test}{\mathcal{D}_{\text{test}}}
\newcommand{\val}{\mathcal{D}_{\text{val}}}

% Data samples
\newcommand{\x}{\vect{x}}
\newcommand{\y}{y}
\newcommand{\X}{\mat{X}}
\newcommand{\Y}{\vect{y}}

% Labels and predictions
\newcommand{\yhat}{\hat{y}}
\newcommand{\Yhat}{\hat{\vect{y}}}
\newcommand{\fhat}{\hat{f}}

% Parameters
\newcommand{\w}{\vect{w}}
\newcommand{\what}{\hat{\vect{w}}}
\newcommand{\wstar}{\vect{w}^{*}}
\newcommand{\theta}{\boldsymbol{\theta}}
\newcommand{\thetahat}{\hat{\boldsymbol{\theta}}}
\newcommand{\thetastar}{\boldsymbol{\theta}^{*}}

% Gradients
\newcommand{\grad}[1]{\nabla_{#1}}
\newcommand{\gradw}{\nabla_{\w}}
\newcommand{\gradtheta}{\nabla_{\theta}}

% Loss functions
\newcommand{\loss}[1]{\mathcal{L}\left(#1\right)}
\newcommand{\losshat}[1]{\hat{\mathcal{L}}\left(#1\right)}

% Regularization
\newcommand{\reg}[1]{\mathcal{R}\left(#1\right)}

% Objective
\newcommand{\obj}[1]{\mathcal{J}\left(#1\right)}

% Learning rate
\newcommand{\lr}{\eta}
\newcommand{\stepsize}{\alpha}

% Activation functions
\DeclareMathOperator{\relu}{ReLU}
\DeclareMathOperator{\sigmoid}{sigmoid}
\DeclareMathOperator{\softmax}{softmax}
\DeclareMathOperator{\tanhop}{tanh}

% Layer notation
\newcommand{\layer}[1]{^{[#1]}}

% =============================================================================
% Deep Learning
% =============================================================================

% Network
\newcommand{\net}{f_{\theta}}
\newcommand{\netn}{f_{\theta}^{(n)}}

% Attention
\DeclareMathOperator{\Attention}{Attention}
\DeclareMathOperator{\MultiHead}{MultiHead}
\DeclareMathOperator{\SelfAttention}{SelfAttention}

% Transformer
\DeclareMathOperator{\LayerNorm}{LayerNorm}
\DeclareMathOperator{\FFN}{FFN}

% Convolution
\DeclareMathOperator{\Conv}{Conv}
\DeclareMathOperator{\Pool}{Pool}

% =============================================================================
% Reinforcement Learning
% =============================================================================

% States and actions
\newcommand{\state}{s}
\newcommand{\action}{a}
\newcommand{\State}{\mathcal{S}}
\newcommand{\Action}{\mathcal{A}}

% Rewards
\newcommand{\reward}{r}
\newcommand{\Reward}{\mathcal{R}}

% Policy
\newcommand{\policy}{\pi}
\newcommand{\policytheta}{\pi_{\theta}}

% Value functions
\newcommand{\vval}{V}
\newcommand{\qval}{Q}
\newcommand{\aval}{A}

% Advantage
\newcommand{\advantage}{A}

% Return
\newcommand{\return}{G}

% Discount factor
\newcommand{\discount}{\gamma}

% =============================================================================
% Common Functions and Operators
% =============================================================================

% Big-O notation
\DeclareMathOperator{\bigO}{\mathcal{O}}
\DeclareMathOperator{\littleo}{o}
\DeclareMathOperator{\bigOmega}{\Omega}
\DeclareMathOperator{\bigTheta}{\Theta}

% Standard functions
\DeclareMathOperator{\softmaxop}{softmax}
\DeclareMathOperator{\logsumexp}{log-sum-exp}

% Indicator function
\newcommand{\indicator}[1]{\mathbb{I}\left[#1\right]}

% Cardinality
\newcommand{\card}[1]{\left|#1\right|}

% Floor and ceiling
\DeclareMathOperator{\floor}{floor}
\DeclareMathOperator{\ceil}{ceil}

% =============================================================================
% Brackets and Delimiters
% =============================================================================

% Parentheses
\newcommand{\paren}[1]{\left(#1\right)}
\newcommand{\bracket}[1]{\left[#1\right]}
\newcommand{\brace}[1]{\left\{#1\right\}}
\newcommand{\absval}[1]{\left|#1\right|}

% Set notation
\newcommand{\set}[1]{\left\{#1\right\}}

% Conditional
\newcommand{\given}{\,|\,}
\newcommand{\suchthat}{\,:\,}

% =============================================================================
% Common Abbreviations
% =============================================================================

% i.i.d.
\newcommand{\iid}{\textit{i.i.d.}}

% w.r.t., w.p., etc.
\newcommand{\wrt}{w.r.t.}
\newcommand{\wp}{w.p.}
\newcommand{\st}{s.t.}

% et al.
\newcommand{\etal}{\textit{et al.}}

% e.g., i.e., cf.
\newcommand{\eg}{e.g.,}
\newcommand{\ie}{i.e.,}
\newcommand{\cf}{cf.}

% =============================================================================
% Custom Operators (add your own below)
% =============================================================================

% Add your custom mathematical notation here
%
% \newcommand{\yourcommand}[n]{definition}
%
% Example:
% \newcommand{\model}{\mathcal{M}}
% \newcommand{\embed}[1]{\text{embed}(#1)}


% =============================================================================
% General Notation
% =============================================================================

% Vectors and matrices
\newcommand{\vect}[1]{\boldsymbol{\mathbf{#1}}}
\newcommand{\mat}[1]{\boldsymbol{\mathbf{#1}}}
\newcommand{\tensor}[1]{\mathcal{#1}}

% Transpose
\newcommand{\T}{^{\top}}
\newcommand{\Tr}{^{\intercal}}

% Inverse
\newcommand{\inv}{^{-1}}

% Hermitian conjugate
\newcommand{\herm}{^{\dagger}}

% Vector and matrix norms
\newcommand{\norm}[1]{\left\|#1\right\|}
\newcommand{\normp}[2]{\left\|#1\right\|_{#2}}
\newcommand{\abs}[1]{\left|#1\right|}

% Trace and determinant
\DeclareMathOperator{\tr}{tr}
\DeclareMathOperator{\Tr}{Tr}
\DeclareMathOperator{\diag}{diag}
\DeclareMathOperator{\Det}{det}

% Sets
\newcommand{\R}{\mathbb{R}}
\newcommand{\N}{\mathbb{N}}
\newcommand{\Z}{\mathbb{Z}}
\newcommand{\C}{\mathbb{C}}

% =============================================================================
% Probability and Statistics
% =============================================================================

% Probability operators
\DeclareMathOperator{\E}{\mathbb{E}}
\DeclareMathOperator{\Var}{Var}
\DeclareMathOperator{\Cov}{Cov}
\DeclareMathOperator{\KL}{KL}
\DeclareMathOperator{\Entropy}{H}
\DeclareMathOperator{\MI}{I}

% Probability commands
\newcommand{\prob}[1]{\mathbb{P}\left(#1\right)}
\newcommand{\expect}[1]{\E\left[#1\right]}
\newcommand{\expectsub}[2]{\E_{#1}\left[#2\right]}
\newcommand{\var}[1]{\Var\left[#1\right]}
\newcommand{\cov}[1]{\Cov\left[#1\right]}

% Distributions
\DeclareMathOperator{\Ndist}{\mathcal{N}}
\DeclareMathOperator{\Unif}{Uniform}
\DeclareMathOperator{\Bern}{Bernoulli}
\DeclareMathOperator{\Cat}{Categorical}

% Gaussian
\newcommand{\gauss}[2]{\Ndist\left(#1, #2\right)}

% =============================================================================
% Optimization
% =============================================================================

% Argmin/argmax
\DeclareMathOperator{\argmin}{arg\,min}
\DeclareMathOperator{\argmax}{arg\,max}
\DeclareMathOperator{\arginf}{arg\,inf}
\DeclareMathOperator{\argsup}{arg\,sup}

% Min/max
\DeclareMathOperator{\minimize}{minimize}
\DeclareMathOperator{\maximize}{maximize}

% Subgradient
\DeclareMathOperator{\sign}{sign}
\DeclareMathOperator{\prox}{prox}

% =============================================================================
% Linear Algebra
% =============================================================================

% Inner products
\newcommand{\inner}[2]{\left\langle #1, #2 \right\rangle}
\newcommand{\dotp}[2]{#1 \cdot #2}

% Outer product
\newcommand{\outerp}[2]{#1 \otimes #2}

% Matrix operations
\DeclareMathOperator{\rank}{rank}
\DeclareMathOperator{\nullity}{nullity}
\DeclareMathOperator{\range}{range}
\DeclareMathOperator{\nullspace}{null}
\DeclareMathOperator{\spanop}{span}

% Eigenvalues
\DeclareMathOperator{\eig}{eig}
\DeclareMathOperator{\eigmax}{\lambda_{max}}
\DeclareMathOperator{\eigmin}{\lambda_{min}}

% Singular values
\DeclareMathOperator{\svd}{svd}
\DeclareMathOperator{\sigmamax}{\sigma_{max}}
\DeclareMathOperator{\sigmamin}{\sigma_{min}}

% =============================================================================
% Machine Learning Notation
% =============================================================================

% Dataset
\newcommand{\dataset}{\mathcal{D}}
\newcommand{\train}{\mathcal{D}_{\text{train}}}
\newcommand{\test}{\mathcal{D}_{\text{test}}}
\newcommand{\val}{\mathcal{D}_{\text{val}}}

% Data samples
\newcommand{\x}{\vect{x}}
\newcommand{\y}{y}
\newcommand{\X}{\mat{X}}
\newcommand{\Y}{\vect{y}}

% Labels and predictions
\newcommand{\yhat}{\hat{y}}
\newcommand{\Yhat}{\hat{\vect{y}}}
\newcommand{\fhat}{\hat{f}}

% Parameters
\newcommand{\w}{\vect{w}}
\newcommand{\what}{\hat{\vect{w}}}
\newcommand{\wstar}{\vect{w}^{*}}
\newcommand{\theta}{\boldsymbol{\theta}}
\newcommand{\thetahat}{\hat{\boldsymbol{\theta}}}
\newcommand{\thetastar}{\boldsymbol{\theta}^{*}}

% Gradients
\newcommand{\grad}[1]{\nabla_{#1}}
\newcommand{\gradw}{\nabla_{\w}}
\newcommand{\gradtheta}{\nabla_{\theta}}

% Loss functions
\newcommand{\loss}[1]{\mathcal{L}\left(#1\right)}
\newcommand{\losshat}[1]{\hat{\mathcal{L}}\left(#1\right)}

% Regularization
\newcommand{\reg}[1]{\mathcal{R}\left(#1\right)}

% Objective
\newcommand{\obj}[1]{\mathcal{J}\left(#1\right)}

% Learning rate
\newcommand{\lr}{\eta}
\newcommand{\stepsize}{\alpha}

% Activation functions
\DeclareMathOperator{\relu}{ReLU}
\DeclareMathOperator{\sigmoid}{sigmoid}
\DeclareMathOperator{\softmax}{softmax}
\DeclareMathOperator{\tanhop}{tanh}

% Layer notation
\newcommand{\layer}[1]{^{[#1]}}

% =============================================================================
% Deep Learning
% =============================================================================

% Network
\newcommand{\net}{f_{\theta}}
\newcommand{\netn}{f_{\theta}^{(n)}}

% Attention
\DeclareMathOperator{\Attention}{Attention}
\DeclareMathOperator{\MultiHead}{MultiHead}
\DeclareMathOperator{\SelfAttention}{SelfAttention}

% Transformer
\DeclareMathOperator{\LayerNorm}{LayerNorm}
\DeclareMathOperator{\FFN}{FFN}

% Convolution
\DeclareMathOperator{\Conv}{Conv}
\DeclareMathOperator{\Pool}{Pool}

% =============================================================================
% Reinforcement Learning
% =============================================================================

% States and actions
\newcommand{\state}{s}
\newcommand{\action}{a}
\newcommand{\State}{\mathcal{S}}
\newcommand{\Action}{\mathcal{A}}

% Rewards
\newcommand{\reward}{r}
\newcommand{\Reward}{\mathcal{R}}

% Policy
\newcommand{\policy}{\pi}
\newcommand{\policytheta}{\pi_{\theta}}

% Value functions
\newcommand{\vval}{V}
\newcommand{\qval}{Q}
\newcommand{\aval}{A}

% Advantage
\newcommand{\advantage}{A}

% Return
\newcommand{\return}{G}

% Discount factor
\newcommand{\discount}{\gamma}

% =============================================================================
% Common Functions and Operators
% =============================================================================

% Big-O notation
\DeclareMathOperator{\bigO}{\mathcal{O}}
\DeclareMathOperator{\littleo}{o}
\DeclareMathOperator{\bigOmega}{\Omega}
\DeclareMathOperator{\bigTheta}{\Theta}

% Standard functions
\DeclareMathOperator{\softmaxop}{softmax}
\DeclareMathOperator{\logsumexp}{log-sum-exp}

% Indicator function
\newcommand{\indicator}[1]{\mathbb{I}\left[#1\right]}

% Cardinality
\newcommand{\card}[1]{\left|#1\right|}

% Floor and ceiling
\DeclareMathOperator{\floor}{floor}
\DeclareMathOperator{\ceil}{ceil}

% =============================================================================
% Brackets and Delimiters
% =============================================================================

% Parentheses
\newcommand{\paren}[1]{\left(#1\right)}
\newcommand{\bracket}[1]{\left[#1\right]}
\newcommand{\brace}[1]{\left\{#1\right\}}
\newcommand{\absval}[1]{\left|#1\right|}

% Set notation
\newcommand{\set}[1]{\left\{#1\right\}}

% Conditional
\newcommand{\given}{\,|\,}
\newcommand{\suchthat}{\,:\,}

% =============================================================================
% Common Abbreviations
% =============================================================================

% i.i.d.
\newcommand{\iid}{\textit{i.i.d.}}

% w.r.t., w.p., etc.
\newcommand{\wrt}{w.r.t.}
\newcommand{\wp}{w.p.}
\newcommand{\st}{s.t.}

% et al.
\newcommand{\etal}{\textit{et al.}}

% e.g., i.e., cf.
\newcommand{\eg}{e.g.,}
\newcommand{\ie}{i.e.,}
\newcommand{\cf}{cf.}

% =============================================================================
% Custom Operators (add your own below)
% =============================================================================

% Add your custom mathematical notation here
%
% \newcommand{\yourcommand}[n]{definition}
%
% Example:
% \newcommand{\model}{\mathcal{M}}
% \newcommand{\embed}[1]{\text{embed}(#1)}
